Soit un graphe $G$ avec $n$ sommets. 
Si il n'y a pas d'arête alors G est bien acyclique et a $n$ composantes connexes.

On ajoute  $k$ arête avec $k<n$ tel que G est toujours acyclique 
On suppose que le graphe a $|V_G| - |E_G|$ composantes connexes.
On ajoute une arête de tel sorte que G soit toujouors acyclique, alors cette arête connecte deux compsantes connexes distinctes. Alors le nombre de composante connexe diminue de un et le nombre d'arête augmente de 1. On a donc toujours $|V_G| - |E_G|$ composantes connexes.


Voici un exemple: 
\includegraphics[width = 0.5\textwidth]{ex2.png}



